\documentclass[11pt]{article}
\usepackage[T1]{fontenc}
\usepackage{lmodern,amsmath,amssymb,amsthm,mathtools,booktabs,microtype}
\usepackage[margin=1.08in]{geometry}
\usepackage{xcolor}
\usepackage[colorlinks=true,linkcolor=blue!45!black,urlcolor=blue!45!black,citecolor=blue!45!black]{hyperref}
\setlength{\emergencystretch}{2em}
\newtheorem{theorem}{Theorem}[section]
\newtheorem{lemma}[theorem]{Lemma}
\newtheorem{corollary}[theorem]{Corollary}
\newtheorem{proposition}[theorem]{Proposition}
\theoremstyle{remark}\newtheorem{remark}[theorem]{Remark}
\newcommand{\E}{\mathbb E}
\newcommand{\U}{\mathsf U}
\newcommand{\Pp}{\mathcal P}
\newcommand{\R}{\mathbb R}
\newcommand{\Z}{\mathbb Z}
\newcommand{\leg}[2]{\left(\frac{#1}{#2}\right)}
\newcommand{\abs}[1]{\left|#1\right|}
\newcommand{\norm}[1]{\left\|#1\right\|}
\newcommand{\support}{\operatorname{supp}}
\title{Exact prime cubature after adversarial erasures\\[3pt]\large Polynomial-height quadratic-residue designs from Dirichlet GRH}
\author{Research note for ProveIt}
\date{27 November 2026}
\begin{document}
\maketitle
\begin{abstract}
Fix $r$ distinct odd primes $q_i\le Y$ and $1\le k\le r$. We construct a rational probability law on primes $p$ in a single dyadic interval for which the signs $(q_i/p)$ are exactly $k$-wise independent. Even after an adversary removes primes of baseline mass at most $\delta$, every remaining prime can be given positive weight between $3/4$ and $5/4$ times its renormalized baseline weight. With $D=\sum_{j\le k}\binom rj$, the required interval starts at $X\gg_k(D+\delta D^2)^2\log^2(2Y)$ for $0\le\delta\le c_k$. A sparse alternative uses at most $D$ primes. The arithmetic input is only Dirichlet GRH; the finite-dimensional argument uses hypercontractivity and the support function of a zonotope. A three-layer example shows that the deletion-versus-bias scale of the finite-dimensional theorem is sharp up to constants for degree one. An exact rational certificate illustrates degree two after a targeted deletion.
\end{abstract}

\section{Statement and scope}
Let $\U$ be the uniform law on $\{-1,1\}^r$, and write $u_A=\prod_{i\in A}u_i$ (including $u_\varnothing=1$). Set
\[
 D=D(r,k)=\sum_{j=0}^k\binom rj,\qquad
 b(t)=\begin{cases}30(t-1)^2(2-t)^2&1<t<2,\\0&\text{otherwise}.\end{cases}
\]
Thus $b\ge0$, $\int b=1$, $\norm b_\infty=15/8$. For integer $X\ge2$, let $\Pp=\{p\text{ prime}:X<p<2X,\ p\notin\{q_1,\dots,q_r\}\}$ and put
\[
 u(p)=(\leg{q_1}{p},\ldots,\leg{q_r}{p}),\quad
 \mu(p)=\frac{b(p/X)}{\sum_{l\in\Pp} b(l/X)}.
\]
For $S\subseteq\Pp$ of positive mass write $\mu_S=\mu|_S/\mu(S)$. Throughout, ``Dirichlet GRH'' means RH for $\zeta$ and the assertion that the nontrivial zeros of every primitive Dirichlet $L$-function have real part $1/2$. The vendored source establishes a stronger statement including imprimitive characters; see Section~\ref{sec:source}.

\begin{theorem}[Prime cubature, using Dirichlet GRH]\label{thm:main}
For each $k\ge1$ there is an effective $C_k>0$ such that, for all $r\ge k$, all distinct odd primes $q_1,\dots,q_r\le Y$, and all
\[
 0\le\delta\le c_k:=\frac{1}{1024\,9^k},\qquad
 X\ge C_k(D+\delta D^2)^2\log^2(2Y),
\]
$\Pp$ is nonempty and the following holds simultaneously for every $S\subseteq\Pp$ with $\mu(\Pp\setminus S)\le\delta$. There are positive rational weights $w_p$ on $S$ such that
\begin{equation}\label{eq:cub}
 \frac34\mu_S(p)\le w_p\le\frac54\mu_S(p),\qquad
 \sum_{p\in S}w_p u_A(p)=\begin{cases}1&A=\varnothing,\\0&1\le |A|\le k.\end{cases}
\end{equation}
In particular, for every $B\subseteq[r]$ with $|B|\le k$ and every $v\in\{-1,1\}^B$,
\begin{equation}\label{eq:indep}
 \sum_{p\in S:\ u_i(p)=v_i\ (i\in B)}w_p=2^{-|B|}.
\end{equation}
The same conclusions hold for $u_i(p)=(p/q_i)$, without altering the order of the cutoff.
\end{theorem}
Here and below the constants in asymptotic notation are independent of $r,Y,X$ unless indexed by $k$. The constant $C_k$ is not numerically evaluated in this note; effectiveness refers to the effective explicit-formula and elementary prime-count estimates used in the proof.

\begin{corollary}[Sparse and high-entropy alternatives]\label{cor:sparse}
Under Theorem~\ref{thm:main}, there exists an exact rational design supported on at most $D$ surviving primes (not necessarily with comparable weights). Every exact positive design has at least $D(r,\lfloor k/2\rfloor)$ distinct sign vectors, hence at least that many primes. The full-support design of \eqref{eq:cub} has
\[
 \max_{p\in S}w_p\ll \frac{\log X}{X},\qquad
 H_\infty(w):=-\log_2\max_p w_p\ge \log_2(X/\log X)-O(1).
\]
\end{corollary}
\begin{corollary}[Deletion measured by cardinality]\label{cor:count}
There is an absolute $K\ge1$ such that the theorem holds for deletion of any $\tau\abs{\Pp}$ primes, with $\tau\le c_k/K$, on replacing $\delta$ in the cutoff by $K\tau$. The set of primes need not be chosen randomly.
\end{corollary}
For fixed $k$, the prime-height bound with no deletions is $O_k(r^{2k}\log^2(2Y))$; for a fixed admissible positive deletion fraction it is $O_k(r^{4k}\log^2(2Y))$. The crossover occurs at mass $\delta\asymp D^{-1}$. These are polynomial in $r$ for each fixed $k$, unlike the exponential height from directly finding primes in every full sign cell. If $k$ is allowed to grow, the constants and $D$ must be retained.

\section{The finite moment theorem}
Let $\Omega$ be a finite set of labeled points, $u:\Omega\to\{-1,1\}^r$ (repeated sign vectors allowed), and let $\mu$ be a positive probability law on it. Define its truncated parity bias by
\[
 \epsilon=\max_{1\le |A|\le\min(4k,r)}\abs{\E_\mu u_A}.
\]
We first prove an unconditional theorem that is useful independently of primes.
\begin{theorem}[Bounded-density repair]\label{thm:repair}
Let $D=D(r,k)$ and $H=9^k+\epsilon D^2$. Suppose that
\begin{equation}\label{eq:eps}
 \epsilon\le\frac{1}{32\,3^kD},\qquad
 \mu(\Omega\setminus S)\le\frac{1}{512H}.
\end{equation}
Then there is a probability law $w$ on $S$ satisfying \eqref{eq:cub} with $p$ replaced by points of $\Omega$. If $\mu$ is rational, $w$ can be chosen rational. If only a target deletion budget $\delta$ is given, sufficient conditions are
\begin{equation}\label{eq:budget}
 \delta\le\frac1{1024\,9^k},\quad
 \epsilon\le\frac1{32\,3^kD},\quad
 \delta\epsilon D^2\le\frac1{1024}.
\end{equation}
\end{theorem}
\begin{lemma}[Three norm bounds]\label{lem:norms}
If $f=\sum_{|A|\le k}a_Au_A$ is real and $\E_\U f^2=1$, then
\begin{align*}
 \abs{\E_\mu f-\E_\U f}&\le\epsilon\sqrt D,\\
 1-\epsilon D\le \E_\mu f^2&\le1+\epsilon D,\\
 \E_\mu f^4&\le 9^k+\epsilon D^2.
\end{align*}
\end{lemma}
\begin{proof}
The Walsh characters are orthonormal for $\U$; thus $\sum_Aa_A^2=1$. Their product is $u_{A\triangle B}$. If $L(f)=\sum_A|a_A|$, then $L(f)\le\sqrt D$ by Cauchy--Schwarz, and $L(fg)\le L(f)L(g)$ by the triangle inequality. For a polynomial $g$ of degree at most $4k$, $|\E_\mu g-\E_\U g|\le\epsilon L(g)$. This proves the first two estimates and bounds the fourth-moment discrepancy by $\epsilon D^2$.

For completeness the required hypercontractive inequality follows by induction on the number of coordinates. In fact
\[
 (\E_\U f^4)^{1/2}\le\sum_A 3^{|A|}a_A^2.
\]
Write $f=g+u_r h$ and let $G,H'$ denote the corresponding weighted sums of squared coefficients of $g,h$. The induction hypothesis and Cauchy--Schwarz give
\[
 \E f^4=\E g^4+6\E g^2h^2+\E h^4
 \le G^2+6GH'+(H')^2\le(G+3H')^2.
\]
The initial case of a constant is equality. Since the degree is at most $k$, this implies $\E_\U f^4\le9^k$. All assertions follow.
\end{proof}
\begin{lemma}[Box criterion]\label{lem:box}
Let $v_x=(u_A(x))_{|A|\le k}\in\R^D$, and let $e_0=(1,0,\ldots,0)$. Given a probability law $\nu$ on $S$ and $0<\eta<1$, there are weights $w_x\in[(1-\eta)\nu(x),(1+\eta)\nu(x)]$ with $\sum w_xv_x=e_0$ if and only if for every real polynomial $f$ of degree at most $k$,
\begin{equation}\label{eq:box}
 \abs{\E_\U f-\E_\nu f}\le\eta\E_\nu|f|.
\end{equation}
\end{lemma}
\begin{proof}
The needed correction $e_0-\sum_x\nu(x)v_x$ must belong to the compact convex set $Z=\{\sum_x t_x\nu(x)v_x:|t_x|\le\eta\}$. In direction $a$, its support function is $\eta\sum_x\nu(x)|a\cdot v_x|$. The separating-hyperplane characterization of a compact convex set, applied to $a$ and $-a$, gives precisely \eqref{eq:box}, with $f=a\cdot v$. This proves both directions. The constant coordinate ensures total weight one.
\end{proof}
\begin{proof}[Proof of Theorem~\ref{thm:repair}]
Write $R=\Omega\setminus S$ and $d=\mu(R)$. It is enough to check \eqref{eq:box} for $\nu=\mu_S$, $\eta=1/4$ and $\norm f_{L^2(\U)}=1$, by homogeneity. Put $g=f-\E_\U f$; then $\norm g_{L^2(\U)}\le1$. By Lemma~\ref{lem:norms} and Cauchy--Schwarz,
\begin{align}
 (1-d)\abs{\E_{\mu_S}f-\E_\U f}
 &=\abs{\E_\mu g-\E_\mu(1_R g)}\notag\\
 &\le\epsilon\sqrt D+\sqrt{2d}.\label{eq:num}
\end{align}
The same lemma and H\"older's inequality give $\E_\mu |f|\ge (\E_\mu f^2)^{3/2}/(\E_\mu f^4)^{1/2}$, because $\E f^2\le(\E|f|)^{2/3}(\E f^4)^{1/3}$. Thus
\begin{equation}\label{eq:den}
 (1-d)\E_{\mu_S}|f|\ge \frac{(1-\epsilon D)^{3/2}}{\sqrt H}-\sqrt{2d}.
\end{equation}
Under \eqref{eq:eps}, $\epsilon D\le1/8$, and
\[
 \epsilon^2 D H=\epsilon^2D\,9^k+\epsilon^3D^3
 \le \frac{1}{1024D}+\frac{1}{32768\,27^k}<\frac1{256}.
\]
Also $\sqrt{2dH}\le1/16$. Hence \eqref{eq:num} is at most $1/(8\sqrt H)$, whereas \eqref{eq:den} is at least $1/(2\sqrt H)$, since $(7/8)^{3/2}-1/16>1/2$. Their ratio is at most $1/4$. Lemma~\ref{lem:box} proves the existence. If the baseline is rational, the feasible region is a nonempty bounded polytope described by rational inequalities and equations, so it has a rational vertex. Finally \eqref{eq:budget} implies $\delta H\le1/512$ and hence \eqref{eq:eps}.
\end{proof}
\begin{remark}
The proof controls signed averages and fourth moments, not total variation from $\U$. The law can be supported on far fewer than $2^r$ points. In particular, the argument does not require a prime in each full sign cell. Positivity on all surviving labels is stronger than merely having the target in their convex hull.
\end{remark}

\section{Arithmetic supply of low-degree moments}\label{sec:arithmetic}
The direct arithmetic section of the ProveIt article contains a uniform smoothed prime-sum lemma (label \texttt{dir:lem:smooth}) and uses low-conductor products of quadratic characters (label \texttt{dir:thm:signs}). Here we adapt that analytic input to an exactly rational baseline and a simultaneous moment repair rather than a single positive-polynomial witness.
\begin{lemma}[Uniform rational bump estimate]\label{lem:bump}
Assume Dirichlet GRH. If $X\ge2$ and $\chi$ is a primitive real character of conductor $f$ (including the zeta case $f=1$), then
\begin{equation}\label{eq:bump}
 \sum_{p}b(p/X)\chi(p)=\frac{1}{\log X}\left(\mathbf1_{f=1} X I_X+O(\sqrt X\log(2f))\right),\quad
 I_X=\int_1^2 b(t)\frac{\log X}{\log(Xt)}\,dt.
\end{equation}
The implied constant is absolute, and $1/2\le I_X\le1$.
\end{lemma}
\begin{proof}
Let $h_X(t)=b(t)\log X/\log(Xt)$ on $(1,2)$ and zero elsewhere. Its derivatives through order three on the interior are uniformly bounded for $X\ge2$; both $h_X$ and its first derivative vanish at the two endpoints. Consequently the Mellin transform $\widehat h_X(s)=\int_1^2 h_X(t)t^{s-1}\,dt$ satisfies
\[
 |\widehat h_X(\sigma+it)|\ll (1+|t|)^{-3}\qquad(-1/2\le\sigma\le2),
\]
by integration by parts in $\log t$ (the third step can include bounded endpoint terms). The constants are uniform in $X$.

Apply the standard explicit formula, equivalently shift the Mellin integral from $\Re s=2$ to $-1/2$ as in \cite{davenport,article}. The integrand is
\[ -(L'/L)(s,\chi)X^s\widehat h_X(s). \] The principal pole contributes $XI_X$ when $f=1$. Nontrivial zeros contribute $-\sum_\rho X^\rho\widehat h_X(\rho)$ with multiplicities. Under Dirichlet GRH, grouping ordinates in unit intervals and using the classical local zero-count bound $O(\log(2f(j+2)))$ shows that the sum is absolutely bounded by $O(\sqrt X\log(2f))$. A possible trivial zero at zero contributes $O(1)$. The functional equation and the absolutely convergent Dirichlet series on the reflected line $\Re(1-s)=3/2$ bound the new vertical integral by $O(X^{-1/2}\log(2f))$. The usual choice of heights avoiding zeros is justified by the same local zero-count bound and the local partial-fraction estimate for $L'/L$; the decay of the Mellin transform makes the horizontal integrals vanish. This is the standard smoothed explicit-formula argument; its constants here depend only on the uniform derivative bounds.

The resulting sum over prime powers differs from its prime part by $O(\sqrt X)$, using Chebyshev's estimate:
\[ \sum_{m\ge2}\vartheta((2X)^{1/m})\ll\sqrt X+X^{1/3}\log X\ll\sqrt X. \] On primes $\log p\,h_X(p/X)=\log X\,b(p/X)$. Dividing by $\log X$ proves \eqref{eq:bump}. The bound on $I_X$ follows from $\log X\le\log(Xt)\le2\log X$.
\end{proof}
\begin{lemma}[Truncated bias of the prime pool]\label{lem:primebias}
For each $k\ge1$ there is an effective $B_k$ such that if $X\ge B_k r^2\log^2(2Y)$, the prime pool is nonempty and
\begin{equation}\label{eq:bias}
 \epsilon\le B_k\frac{\log(2Y)}{\sqrt X},\qquad
 Z_X:=\sum_{p\in\Pp}b(p/X)\gg X/\log X,\qquad |\Pp|\ll X/\log X.
\end{equation}
The last two constants can be taken absolute after increasing the lower bound on $X$ by an absolute factor.
\end{lemma}
\begin{proof}
The quadratic character $p\mapsto(q_i/p)$ for odd $p$ not dividing $q_i$ is induced by the primitive Kronecker character of fundamental discriminant $q_i$ if $q_i\equiv1\pmod4$, and $4q_i$ if $q_i\equiv3\pmod4$. For a nonempty subset $A$ the product has primitive character associated to the squarefree positive integer $\prod_{i\in A}q_i$ and its fundamental discriminant. The conductor $f_A$ is either that product or four times it, and the character is nonprincipal. For $(p/q_i)$ the CRT gives the same conclusion with conductor $\prod_{i\in A}q_i$. For $|A|\le4k$, $\log(2f_A)\ll k\log(2Y)$. The empty subset uses $f=1$.

Apply Lemma~\ref{lem:bump} to these primitive characters. Excluding all $q_i$ costs at most $r\norm b_\infty$ in an unweighted sum. More precisely, before division by $\log X$ it costs $O(r\log(2Y))$: if an excluded prime lies in $(X,2X)$ then $X<q_i\le Y$. Thus the principal sum has main term $XI_X$ and error $O(\sqrt X+r\log(2Y))$, and every nonempty parity sum has size $O_k(\sqrt X\log(2Y)+r\log(2Y))$, both before division by $\log X$. The bound on $X$ makes the principal sum at least $X/4$ at this scale and absorbs the exclusion term in the nonprincipal bound. Dividing the latter by the principal sum proves the first two assertions of \eqref{eq:bias}. Chebyshev's prime-count bound gives the third.
\end{proof}
\begin{proof}[Proof of Theorem~\ref{thm:main}]
Because $D\ge r+1$, increasing $C_k$ gives the threshold in Lemma~\ref{lem:primebias}. We can further increase $C_k$ so that \eqref{eq:bias} implies
\[
 \epsilon\le \frac{1}{32\,3^kD+1024\delta D^2}.
\]
This ensures all conditions \eqref{eq:budget}, including when $\delta=0$. Apply the unconditional repair theorem. The baseline is rational for integer $X$, and every prime in the open interval has positive baseline weight. Formula \eqref{eq:indep} follows by expanding the cell indicator $2^{-|B|}\prod_{i\in B}(1+v_i u_i)$. No character of order higher than two and no nonabelian $L$-function is used.
\end{proof}
\begin{proof}[Proof of Corollaries~\ref{cor:sparse} and \ref{cor:count}]
The exact weights solve $\sum_p w_p v_p=e_0$ with $w_p\ge0$ and rational coefficients. The feasible polytope has a rational extreme point supported on at most $D$ columns: if there were more, a nonzero linear dependence on its support would allow a small perturbation in either direction. For the lower bound, the evaluation vectors of Walsh polynomials of degree at most $\lfloor k/2\rfloor$ must be orthonormal in the weighted inner product, because their pairwise products have degree at most $k$. Their number cannot exceed the number of distinct sign vectors in the support.

By \eqref{eq:bias}, $\mu(p)\ll \log X/X$; renormalizing after deletion multiplies it by at most $(1-c_k)^{-1}$, and \eqref{eq:cub} by at most $5/4$. This proves the entropy bound. Moreover, if $R$ has at most $\tau|\Pp|$ primes, its baseline mass is at most $\tau|\Pp|\norm b_\infty/Z_X\le K\tau$ for an absolute $K$. This proves the cardinality version.
\end{proof}

\section{A sharp obstruction for linear tests}
The following example shows that the $\epsilon D^2$ term in the permitted deletion scale of Theorem~\ref{thm:repair} has real content. It concerns arbitrary finite probability laws on the cube, not a lower bound for primes.
\begin{theorem}[Three-layer obstruction]\label{thm:sharp}
Let $r\ge7$ be odd, and choose an odd integer $a\le r$ with $a^2\ge6r$. There is a rational exchangeable law $\mu$ on the cube such that its one- and two-coordinate parity moments vanish, its maximum bias through degree four is
\[
 \epsilon\asymp a^2/r^3,
\]
and deleting a set of mass $d\asymp r/a^2$ leaves a support admitting no exact one-wise independent probability law, irrespective of its weights. In particular $d\epsilon r^2\asymp1$. The constants in these comparisons are absolute.
\end{theorem}
\begin{proof}
Write $T=\sum_i u_i$. Let $\mu$ be uniform on each of the three layers $T=1,a,-a$ with respective masses $1-\lambda,\alpha,\beta$, where
\[
 \lambda=\frac{r-1}{a^2-1},\qquad
 \alpha=\frac12\left(\lambda-\frac{1-\lambda}{a}\right),\quad
 \beta=\frac12\left(\lambda+\frac{1-\lambda}{a}\right).
\]
All masses are nonnegative: $\lambda\ge(1-\lambda)/a$ is equivalent to $a\le r$. The choices give $\E T=0$, $\E T^2=r$, $\E T^3=r-a^2$ and $\E T^4=r+(r-1)a^2$. For $j$ distinct coordinates, let $m_j$ be their common parity moment. Expanding powers of $T$, or using the elementary symmetric polynomials of the signs, gives
\begin{align*}
 m_1&=\frac{\E T}{r},\qquad
 m_2=\frac{\E T^2-r}{r(r-1)},\\
 m_3&=\frac{\E T^3-(3r-2)\E T}{r(r-1)(r-2)},\\
 m_4&=\frac{\E T^4-(6r-8)\E T^2+3r(r-2)}{r(r-1)(r-2)(r-3)}.
\end{align*}
These formulas follow also from the coefficient of $z^j$ in $(1+z)^{(r+T)/2}(1-z)^{(r-T)/2}$ divided by $\binom rj$. Substitution yields
\[
 m_1=m_2=0,\quad m_3=\frac{r-a^2}{r(r-1)(r-2)},\quad
 m_4=\frac{a^2-3r}{r(r-2)(r-3)}.
\]
Since $6r\le a^2\le r^2$, the claimed bias comparison follows. Delete the layer $T=-a$, of mass $d=\beta\in[\lambda/2,\lambda]$. The remaining layers satisfy $T\ge1$, so any probability law on them has $\E T>0$. Exact first moments would imply $\E T=0$, impossible. The mass comparison and product estimate follow from the expression for $\lambda$.
\end{proof}
As $a$ varies from order $\sqrt r$ to $r$, the obstruction ranges from a constant deletion fraction at bias of order $r^{-2}$ to deletion of order $r^{-1}$ at bias of order $r^{-1}$. In the intermediate region the necessary relationship for a guarantee using only degree-four biases is, up to constants, $d\lesssim1/(\epsilon r^2)$. This matches the scale of Theorem~\ref{thm:repair} for $k=1$, with $D=r+1$. It does not prove that the prime-height exponent is sharp or that higher-degree bias cannot improve the guarantee.

\section{Construction and reproducible certificate}
For any concrete pool, form the matrix $V$ with rows $v_p=(u_A(p))_{|A|\le k}$ and baseline vector $\nu=\mu_S$. The general algorithm is a rational linear feasibility problem with the boxes in \eqref{eq:cub}; an exact LP solver, or enumeration of basic feasible solutions, suffices. A useful faster special construction, not guaranteed to meet the boxes in every instance, is
\[
 G=V^{\mathsf T}\operatorname{diag}(\nu)V,\qquad
 c=G^{-1}e_0,\qquad w_p=\nu(p)\,(Vc)_p.
\]
When $G$ is invertible and $3/4\le(Vc)_p\le5/4$, this gives an exact rational certificate. Otherwise use the box feasibility problem, not an unverified numerical approximation.

In the supplied example $r=8$, $k=2$, $q=(3,5,7,11,13,17,19,23)$ and $X=100000$. Delete primes with the first two signs $+1$, taking largest bump weights first until the deleted mass reaches the stated budget. The generator uses integer unnormalized bump weights $(p-X)^2(2X-p)^2$, hence does not need logarithms. The independent verifier enumerates primes by trial division, computes symbols using Euler's criterion, and verifies every equality and bound using rational arithmetic.
\begin{center}
\begin{tabular}{lr}\toprule
Prime pool size&8392\\
Deleted primes&94\\
Surviving primes&8298\\
Exact moment equations&37\\
Deleted baseline mass (rounded)&0.02000000\\
Minimum weight ratio (rounded)&0.87583660\\
Maximum weight ratio (rounded)&1.19291444\\
Three-layer moment tests&906\\\bottomrule
\end{tabular}
\end{center}
These are finite certificates, not estimates of $C_k$ or a verification of the general analytic theorem. The deletion in the example exceeds the conservative universal $c_2$; this particular deletion is independently certified and no universal assertion at that size is inferred.

\section{Relation to the source and earlier work}\label{sec:source}
\paragraph{Formal input.} The vendor README \cite{vendor} identifies the imported readable edition of \texttt{lsernau/grh-lean} at commit \texttt{72fead102f663d5c7a407964a13dc0cf9d977207}. It pins Lean v4.35.0 and Mathlib commit \texttt{56a005855ea17e0b0958f73ec801ae140b0a1679}. We read \texttt{src/Statements.lean}, which states \texttt{DirichletLZerosOnLine} and \texttt{DirichletLNonvanishing}\allowbreak\texttt{RightOfHalf} using Mathlib's definitions, and \texttt{src/Main.lean}, which proves them via \texttt{New.zeros\_on\_line New.grh} and \texttt{New.grh}. The README reports axiom dependencies \texttt{propext}, \texttt{Classical.choice}, \texttt{Quot.sound}. The independent verification account \cite{waterloo} reports successful checking of the compact edition of the same statements; the Lean community statement \cite{leanblog} certifies RH. We did not replay either kernel here. The vendor is a separate Lake package and is not checked by ProveIt's root build. The supplied Lean audit file is a reproducible entry point, not a formalization of this paper.

\paragraph{Article comparison.} The article README describes a merge of 131 reports, 723 pages and 914 results, with no independent review of all merged mathematics. Its source section (\texttt{src:thm:dgrh}) identifies the input theorem; its introduction describes the comparison, spanning, twisting, and smoothed prime-sum mechanism. The direct arithmetic section contains (i) the uniform smoothed prime-sum bound \texttt{dir:lem:smooth}, (ii) prescribed-majority sign witnesses \texttt{dir:thm:signs}, and (iii) resilient squarefree product mixing \texttt{dir:thm:quotient}. Our application uses the same smoothing and character parity principles, but asks for a single exact law integrating \emph{all} degree-$k$ tests simultaneously, with bounded changes of weights on every surviving prime and a bias/deletion tradeoff. This differs from finding one prime in a Hamming cap or multiplying many distinct primes to hit a residue class. We inspected the article's introduction, table of sections and the relevant direct arithmetic material, not all 914 proofs. No theorem in those inspected sections was found with the present formulation.

\paragraph{Literature and novelty.} Small-bias spaces and almost $k$-wise independence are classical \cite{naor}. Prescribing Legendre signs of integers is studied by Hanson, Vaughan and Zhang \cite{hvz}. Fourier complexity in effective Chebotarev is developed by Bella\"iche \cite{bellaiche}; only its abelian Dirichlet setting is relevant here. Exact positive cubature and stable randomized constructions are studied in \cite{glaubitz,migliorati}. The box criterion and hypercontractivity are standard tools. The proposed contribution is their specific combination: a quantitative deterministic prime-height bound for exact bounded-density moment repair under arbitrary erasures, and the matching linear-case bias obstruction. A targeted literature search found the above related work but does not establish priority. No claim is made that this is a previously unpublished theorem in an exhaustive sense.

\paragraph{Further directions.} Improve the height exponent for fixed positive deletion fractions; find a sharp higher-degree obstruction; replace bounded reweighting with equal weights on distinct primes; evaluate effective constants; formalize the finite theorem and then the smoothed explicit formula in Lean. The first and third are not established here.

\begin{thebibliography}{10}\small
\bibitem{article} ProveIt Contributors, \emph{Beyond the Riemann hypothesis for Dirichlet L-functions: consequences, transfers and limits}, 26 Nov. 2026, \href{https://github.com/VladimirReshetnikov/ProveIt/tree/main/Analysis/GeneralizedRiemannHypothesis/Article}{repository source}.
\bibitem{vendor} L. Sernau et al., \emph{The generalized Riemann hypothesis for Dirichlet L-functions, in Lean 4}, readable edition, \href{https://github.com/lsernau/grh-lean/tree/72fead102f663d5c7a407964a13dc0cf9d977207/grh}{pinned upstream source}.
\bibitem{waterloo} M. Rubinstein, \emph{Alp\"oge/Claude proof of the generalized Riemann hypothesis}, verification account, \url{https://www.math.uwaterloo.ca/~mrubinst/alpoge_claude_grh/}, accessed 27 Nov. 2026.
\bibitem{leanblog} Lean FRO and Mathlib Initiative, \emph{Lean verification of the Riemann hypothesis}, \url{https://leanprover-community.github.io/blog/posts/riemann-hypothesis/}, accessed 27 Nov. 2026.
\bibitem{davenport} H. Davenport, \emph{Multiplicative Number Theory}, 3rd ed., Springer, 2000 (Dirichlet $L$-functions and explicit formulas).
\bibitem{naor} J. Naor and M. Naor, Small-bias probability spaces: efficient constructions and applications, \emph{SIAM J. Comput.} 22 (1993), 838--856. \href{https://doi.org/10.1137/0222053}{doi:10.1137/0222053}.
\bibitem{hvz} B. Hanson, R. C. Vaughan and R. Zhang, The least number with prescribed Legendre symbols and representation by binary quadratic forms of small discriminant, \emph{J. Number Theory} 179 (2017), 3--16. \href{https://doi.org/10.1016/j.jnt.2017.03.004}{doi:10.1016/j.jnt.2017.03.004}.
\bibitem{bellaiche} J. Bella\"iche, Th\'eor\`eme de Chebotarev et complexit\'e de Littlewood, \emph{Ann. Sci. ENS} 49 (2016), 579--632. \href{https://doi.org/10.24033/asens.2291}{doi:10.24033/asens.2291}.
\bibitem{glaubitz} J. Glaubitz, Construction and application of provable positive and exact cubature formulas, \emph{IMA J. Numer. Anal.} 43 (2023), 1616--1652. \href{https://academic.oup.com/imajna/article/43/3/1616/6586010}{Article}.
\bibitem{migliorati} G. Migliorati and F. Nobile, Stable high-order randomized cubature formulae in arbitrary dimension, \emph{J. Approx. Theory} 275 (2022), 105706. \href{https://doi.org/10.1016/j.jat.2022.105706}{doi:10.1016/j.jat.2022.105706}.
\end{thebibliography}
\end{document}
